\documentclass[a4paper]{article}

\usepackage[UTF8]{ctex}
\usepackage[a4paper,margin=1in]{geometry}
\usepackage{graphicx}
\usepackage{float}
\usepackage{listings}
\usepackage{booktabs}
\usepackage{color}
\usepackage{indentfirst}
\setlength{\parindent}{2em}
\usepackage{zhnumber}
\usepackage{adjustbox}
\usepackage{subcaption}
\usepackage{amsmath}
\usepackage{amssymb}
\usepackage{enumitem}
\usepackage{lmodern}
\usepackage{inconsolata}
\usepackage{multirow}
\usepackage{tikz}
\usetikzlibrary{positioning}
\usepackage{fancyhdr}
\usepackage{hyperref}

% 配置lstlisting支持中文
\lstset{
    language=C,
    basicstyle=\footnotesize\ttfamily,
    breaklines=true,
    keywordstyle=\bfseries\color{blue},
    commentstyle=\itshape\color{green!50!black},
    stringstyle=\bfseries\color{red},
    numbers=left,
    numberstyle=\tiny\color{gray},
    frame=single,
    backgroundcolor=\color{gray!5},
    showstringspaces=false,
    keepspaces=true,
    tabsize=4,
    extendedchars=true,
    inputencoding=utf8,
    literate={，}{{\textcolor{green!50!black}{，}}}{1}
             {。}{{\textcolor{green!50!black}{。}}}{1}
             {：}{{\textcolor{green!50!black}{：}}}{1}
             {；}{{\textcolor{green!50!black}{；}}}{1}
             {（}{{\textcolor{green!50!black}{（}}}{1}
             {）}{{\textcolor{green!50!black}{）}}}{1}
             {【}{{\textcolor{green!50!black}{【}}}{1}
             {】}{{\textcolor{green!50!black}{】}}}{1}
}

\newcommand{\projname}{计算机网络实验|实验六}
\newcommand{\reporttitle}{VLAN个人心得体会}
\newcommand{\groupno}{第一组}
\newcommand{\authorname}{马福泉}
\newcommand{\major}{计算机科学与技术}
\newcommand{\adviser}{谢逸}
\newcommand{\startdate}{2025年11月8日}
\newcommand{\labenddate}{2025年11月22日}
\newcommand{\labroom}{实验D402}

\pagestyle{fancy}
\fancyhf{}
\fancyhead[L]{\kaishu \projname}
\fancyhead[C]{\kaishu \reporttitle}
\fancyhead[R]{\kaishu \authorname}
\fancyfoot[C]{第 \thepage 页}
\renewcommand{\headrulewidth}{0.4pt}
\renewcommand{\footrulewidth}{0pt}


\begin{document}

% 封面
\begin{titlepage}
    \centering
    \includegraphics[width=10cm]{img/SYSULogo.png}
    \vspace{1em}
    % 修改为换行并居中
    \par
    {\Large \kaishu
        \par
        \projname
        \par
    }
    \vspace{3em}
    {\fontsize{40pt}{42pt}\kaishu \selectfont \boldmath \reporttitle\par}
    \vspace*{\fill}
    \begin{center}
    {\Large
    \makebox[5em][s]{姓名}:\underline{\makebox[15em][c]{\kaishu 马福泉}}\\[1em]
    \makebox[5em][s]{学号}:\underline{\makebox[15em][c]{\kaishu 23336179}}\\[1em]
    \makebox[5em][s]{小组编号}:\underline{\makebox[15em][c]{\kaishu \groupno}}\\[1em]
    \makebox[5em][s]{专业}:\underline{\makebox[15em][c]{\kaishu \major}}\\[1em]
    \makebox[5em][s]{课程教师}:\underline{\makebox[15em][c]{\kaishu \adviser}}\\[1em]
    \makebox[5em][s]{起始日期}:\underline{\makebox[15em][c]{\kaishu \startdate}}\\[1em]
    \makebox[5em][s]{结束日期}:\underline{\makebox[15em][c]{\kaishu \labenddate}}\\[1em]
    \makebox[5em][s]{实验地点}:\underline{\makebox[15em][c]{\kaishu \labroom}}
    }
    \end{center}
    \vspace*{\fill}
\end{titlepage}

\section{承担的工作}
在本次VLAN跨交换机实现实验中，我主要承担了以下工作：
\subsection{前期拓扑图设计}
\begin{itemize}
    \item \textbf{网络拓扑规划}：根据实验要求，我和队友共同设计了基础实验和拓展实验的网络拓扑结构
    \begin{itemize}
        \item 基础实验：设计PC1、PC2、PC3与交换机A、交换机B的连接关系
        \item 拓展实验：设计分层网络架构，包含弱电箱区域（S1）和客厅区域（S2）
    \end{itemize}
    \item \textbf{VLAN规划设计}：制定基于业务类型的VLAN划分方案
    \begin{itemize}
        \item VLAN 10：普通用户业务（PC1、PC2、网络源L1）
        \item VLAN 20：高级用户业务（PC3）
        \item VLAN 30：IPTV专用业务（IPTV源L2、机顶盒L3）
    \end{itemize}
    \item \textbf{物理连接设计}：确定各设备间的物理连线方式和端口分配方案
\end{itemize}

\subsection{基础实验配置}
\begin{itemize}
    \item \textbf{实验环境搭建}：
    \begin{itemize}
        \item 负责连接PC1、PC2、PC3与交换机A、交换机B的物理连线
        \item 按照设计的拓扑图完成所有设备的物理连接
        \item 验证物理连接的正确性和稳定性
    \end{itemize}
    \item \textbf{主机IP配置}：
    \begin{itemize}
        \item 本人负责为PC3主机配置静态IP地址（192.168.10.30）
        \item 使用netsh命令进行网络接口配置
        \item 进行初始连通性测试，验证所有设备在同一广播域内可以正常通信
    \end{itemize}
\end{itemize}

\subsection{实现与调试}
本人协助队友共同实现了整个实验，全程参与，下面仅列出主要负责部分
\begin{itemize}
    \item \textbf{Trunk链路配置}：
    \begin{itemize}
        \item 在交换机A的G0/0/24端口配置Trunk模式
        \item 在交换机B的G0/0/24端口配置Trunk模式
        \item 验证Trunk链路的正常工作状态
    \end{itemize}
    \item \textbf{拓展实验高级配置}：
    \begin{itemize}
        \item 在S1交换机上创建VLAN 10、20、30并进行端口划分
        \item 实现ACL访问控制，精确控制不同VLAN间的访问权限
    \end{itemize}
    \item \textbf{连通性测试}：
    \begin{itemize}
        \item 系统性地测试各设备间的网络连通性
        \item 验证同VLAN内设备可以正常通信
        \item 验证不同VLAN间设备无法直接通信
        \item 测试跨交换机的VLAN通信功能
    \end{itemize}
    \item \textbf{业务功能验证}：
    \begin{itemize}
        \item 验证PC1、PC2、PC3均能访问互联网（L1）
    \end{itemize}
    \item \textbf{故障排除}：
    \begin{itemize}
        \item 排查Trunk链路配置问题，发现交换机B未创建相应VLAN
        \item 解决Windows防火墙阻止ICMP协议的问题
        \item 调试ACL规则，确保访问控制的精确性
        \item 使用display命令检查配置状态，定位问题原因
    \end{itemize}
\end{itemize}

\subsection{实验报告编写}
\begin{itemize}
    \item \textbf{文档整理与格式化}：
    \begin{itemize}
        \item 统一代码块格式和图片布局
        \item 完善目录结构和章节编号
        \item 检查文档的完整性和准确性
    \end{itemize}
\end{itemize}

\section{遇到的困难及解决方法}

\subsection{Trunk链路配置问题}
\textbf{困难描述：}在配置交换机间Trunk链路时，虽然配置了\lstinline{port trunk allow-pass vlan all}，但PC2和PC3仍无法通信。

\textbf{解决方法：}
\begin{itemize}
    \item 发现问题在于交换机B上未创建相应的VLAN，导致VLAN 20的数据无法正确处理
    \item 在交换机B上创建VLAN 20，并将PC3所在端口划分到VLAN 20
    \item 验证两端交换机都配置了相同的Trunk模式，确保VLAN数据能够双向传输
\end{itemize}

\subsection{拓展实验中的ACL配置挑战}
\textbf{困难描述：}在实现复杂的访问控制需求时，ACL规则配置比较复杂，需要精确控制不同VLAN间的访问权限。

\textbf{解决方法：}
\begin{itemize}
    \item 仔细分析业务需求，明确哪些设备间需要通信，哪些需要隔离
    \item 使用高级ACL（编号3025）创建精细化访问控制规则
\end{itemize}

\subsection{防火墙干扰问题}
\textbf{困难描述：}在测试IPTV设备（L2和L3）连通性时，ping测试失败，怀疑是网络配置问题。

\textbf{解决方法：}
\begin{itemize}
    \item 经过排查发现是Windows防火墙阻止了ICMP协议
    \item 关闭L2和L3设备的防火墙后，ping测试成功
    \item 学会了在网络测试中要考虑操作系统层面的安全设置影响
\end{itemize}

\section{体会与总结}

\subsection{技术技能提升}
通过本次实验，我在网络技术方面获得了显著提升：
\begin{itemize}
    \item \textbf{VLAN技术掌握}：深入理解了VLAN的工作原理，掌握了基于端口的VLAN划分方法
    \item \textbf{交换机配置能力}：熟练掌握了华为交换机的基本配置命令，能够独立完成VLAN创建、端口划分、Trunk配置等操作
    \item \textbf{网络规划设计}：学会了根据实际业务需求进行VLAN规划，实现网络的逻辑分段和流量隔离
    \item \textbf{故障排除能力}：提高了网络故障诊断和解决能力，学会使用各种显示命令检查配置状态
\end{itemize}

\subsection{实验方法论收获}
\begin{itemize}
    \item \textbf{系统性思维}：学会了将复杂的网络需求分解为具体的配置步骤，按部就班地实现目标
    \item \textbf{测试驱动}：每完成一个配置步骤都要进行连通性测试，及时发现和解决问题
    \item \textbf{文档化习惯}：养成了记录配置命令和测试结果的习惯，便于问题追溯和经验总结
\end{itemize}

\subsection{团队协作体验}
在团队协作过程中，我学会了：
\begin{itemize}
    \item \textbf{任务分工}：根据各自的技术特长合理分配实验任务，提高实验效率
    \item \textbf{知识共享}：在遇到困难时主动与队友交流讨论，共同寻找解决方案
    \item \textbf{相互支持}：在队友遇到问题时提供帮助，形成了良好的团队学习氛围
\end{itemize}

总的来说，本次VLAN实验让我从理论走向实践，不仅掌握了具体的技术技能，更重要的是培养了系统性的网络思维和解决问题的能力。这为我后续的专业学习奠定了坚实的基础。
\end{document}
